% TOP_PROBLEM: {"id": 1506, "title": "Class Number Problem", "category_id": 1, "category": {"id": 1, "name": "number_theory", "display_name": "Number Theory", "description": "Properties of integers, prime numbers, Diophantine equations.", "slug": "number-theory", "order_index": 1, "created_at": "2026-07-31T15:26:25.670Z"}, "status": "open"}
% FIRST_POSTED: null
% ENTRY_KIND: statement_only
% Imported from: batch08.tex; UnsolvedMath ID 1506
\documentclass[11pt]{article}
\usepackage[margin=25mm]{geometry}
\usepackage{fontspec}
\setmainfont{Latin Modern Roman}
\usepackage{amsmath,amssymb,mathtools}
\usepackage{unicode-math}
\setmathfont{Latin Modern Math}
\usepackage{xurl,hyperref}
\hypersetup{colorlinks=true,urlcolor=blue}
\setcounter{secnumdepth}{0}
\setlength{\parindent}{0pt}
\setlength{\parskip}{5pt}
\setlength{\emergencystretch}{3em}
\newcommand{\sref}[2]{\hyperlink{source:#1:#2}{[S#2]}}
\newcommand{\cataloguescope}[1]{\paragraph{Recorded target.}#1}
\begin{document}
% UnsolvedMath ID 1506; source catalogue code NUM-012
\setcounter{section}{1505}
\section{Class Number Problem}\label{1506}
{\small\sffamily UnsolvedMath ID 1506 · Number Theory}
\subsection{Definitions and mathematical statement}

\paragraph{Shared notation.}$\mathbb N=\{1,2,\ldots\}$, and $\mathbb Z,\mathbb Q,\mathbb R,\mathbb C$ denote the integers, rationals, reals and complex numbers. Any other conventions are specified locally.


Let \(d>1\) be a square-free integer, meaning no square of a
prime divides \(d\), and let \(K_d=\mathbb Q(\sqrt d)\) be its
real quadratic field. Its ring of algebraic integers is
\[
 \mathcal O_{K_d}=\begin{cases}
   \mathbb Z[(1+\sqrt d)/2],&d\equiv1\pmod4,\\
   \mathbb Z[\sqrt d],&d\equiv2\text{ or }3\pmod4.
 \end{cases}
\]
The ideal class group \(\operatorname{Cl}(K_d)\) is the group
of nonzero fractional ideals of \(\mathcal O_{K_d}\), modulo
the subgroup of principal ideals. Its order \(h(K_d)\) is the class number.
The condition \(h(K_d)=1\) is equivalent to every nonzero ideal
being principal and to unique factorization in \(\mathcal O_{K_d}\).
\paragraph{Statement recorded in the supplied source.}
Are there infinitely many square-free integers \(d>1\)
such that
\[
                          h(\mathbb Q(\sqrt d))=1?
\] \sref{1506}{2}
\subsection{Short English statement}
Do infinitely many different real quadratic number fields have a ring of integers with unique factorization?
\subsection{Sources}
\begin{description}
\item[\hypertarget{source:1506:1}{S1}] ULAM.AI and the UnsolvedMath Contributors, \emph{UnsolvedMath: A Curated Collection of Open Mathematics Problems}, record 1506 (NUM-012). CC BY 4.0. \url{https://huggingface.co/datasets/ulamai/UnsolvedMath}.
\item[\hypertarget{source:1506:2}{S2}] Azizul Hoque and Srinivas Kotyada, \emph{Class number one problem for the real quadratic fields $\mathbb Q(\sqrt{m^2+2r})$}, arXiv:2008.03505 (2020), Section 1, Gauss's unrestricted real-quadratic class-number-one question. \url{https://arxiv.org/abs/2008.03505}.
\end{description}
\label{end:1506}
\end{document}
